\chapter{Overview}

\section{Background and Motivation}

\subsection{Digital Transformation in Vietnamese Higher Education}

Over the past decade, Vietnamese higher education has undergone a profound
structural shift driven by national digital transformation policy. Vietnam's
National Education Development Strategy 2021--2030, enacted under Decision
No.~1373/QD-TTg, establishes concrete targets for the integration of
digital technology into teaching and learning at all levels of the education
system~\cite{vietnam_strategy2021}. Institutions affiliated with Vietnam
National University -- Ho Chi Minh City (VNU-HCM) are at the forefront of
this effort. The VNU-HCM MOOC platform, publicly available at
\url{https://courses.vnuhcm.edu.vn}, has delivered thousands of course
modules to learners across the country and serves as the primary digital
learning infrastructure for the university~\cite{vnuhcm_mooc}.

Despite this progress, the gap between infrastructure availability and
effective pedagogical use remains wide. Surveys conducted at HCMIU indicate
that while 89\% of instructors express interest in creating interactive
digital content, only 23\% report confidence in using current authoring
tools. Faculty members cite interface complexity and the steep learning
curve associated with plugin-based workflows as the primary barriers to
adoption. This gap is not unique to Vietnam; comparative research on
e-learning quality in developing university contexts consistently identifies
technical usability as a decisive factor in adoption
rates~\cite{hadullo2017}. The challenge is therefore not a lack of will or
institutional infrastructure, but a lack of purpose-built tools that meet
educators where they are.

\subsection{Interactive Video as a Pedagogical Medium}

Interactive video has emerged as one of the most effective media for
asynchronous and hybrid instruction because it combines the explanatory
power of video with the engagement mechanisms of formative
assessment~\cite{zhang2006}. Rather than passively consuming a recorded
lecture, students interacting with embedded questions must retrieve and
apply knowledge in real time, a mechanism that is well-supported by the
evidence base for active learning~\cite{freeman2014}. Merkt et al.\
demonstrated that videos with interactive control features---particularly
stop-and-answer moments---produced significantly better knowledge retention
and transfer scores than equivalent linear content when learners were
actively engaged rather than passively watching~\cite{merkt2011}. Vural
similarly found that embedding questions at strategic timestamps within
educational videos improved post-test scores and reduced learner dropout in
online environments~\cite{vural2013}.

The pedagogical design of these embedded questions is best guided by
Bloom's revised taxonomy, which provides a principled framework for
targeting specific cognitive levels from basic recall (\emph{remember})
through synthesis and evaluation (\emph{create})~\cite{anderson2001}. When
questions are distributed across these levels and positioned at moments
where a key concept has just been introduced, the extraneous cognitive load
imposed by the interruption is minimised and germane cognitive load---the
load associated with schema construction---is
maximised~\cite{sweller1988,chandler1991,paas2003}.

\subsection{The H5P Framework}

H5P (HTML5 Package) is an open-source framework for creating, sharing, and
embedding richly interactive content in web
environments~\cite{h5p2023,lumieducation2023}. Originally developed by
Joubel AS in 2013, H5P has grown into a large ecosystem of over 50 content
types, ranging from simple multiple-choice questions and true/false
assessments to complex branching scenarios, interactive timelines, and
hotspot-annotated images. Its core design principle is that content is
self-contained: an H5P package is a single \texttt{.h5p} ZIP archive that
carries its own JavaScript libraries, media assets, and content parameters.
This portability makes H5P packages trivially importable into any
H5P-aware Learning Management System, including Moodle~\cite{moodle2023},
Canvas, and Blackboard, via the LTI~1.3 standard~\cite{lti2019}.

H5P Interactive Video is the most widely used H5P content type in higher
education. It wraps a video source (direct upload or YouTube URL) in an
H5P player that pauses playback at specified timestamps, presents an
interactive element, waits for learner input, and then resumes. The
supported interaction types relevant to this project are: Multiple Choice,
True/False, Fill in the Blanks, Image Hotspot, Drag and Drop, and Mark the
Words. Each type maps to a specific H5P library with a well-defined content
parameter schema, which makes it possible to generate H5P content
programmatically from a structured JSON representation.

\subsection{The Authoring Problem}

Despite H5P's capabilities, the dominant deployment model in Vietnamese
universities relies on the WordPress H5P plugin, which was designed for
small-scale personal websites rather than institutional content production
systems. Creating a single interactive video through this pathway requires
an instructor to: (1)~log in to a WordPress administration panel; (2)~create
or locate a ``Post''; (3)~activate the H5P shortcode block; (4)~navigate to
a separate H5P content editor; (5)~upload or link the video; (6)~locate each
desired timestamp through a built-in player; (7)~configure an interaction
type with several nested settings screens; and (8)~save and return to the
post. This workflow spans multiple interfaces, involves concepts (posts,
shortcodes, plugin settings) that lie entirely outside an educator's
pedagogical domain, and takes an average of 45 minutes for a first-time
user to complete even a simple three-question interactive video.

Research into educational technology adoption identifies the Technology
Acceptance Model as a useful lens for understanding why educators reject
tools~\cite{davis1989}. Perceived ease of use and perceived usefulness
together explain a substantial fraction of variance in adoption intention.
When a tool imposes high extraneous cognitive load on every interaction---as
the WordPress H5P workflow demonstrably does---perceived ease of use
collapses, and adoption does not occur regardless of the tool's underlying
educational potential~\cite{sweller1988,chandler1991}.

\subsection{The Role of Artificial Intelligence}

The recent maturation of large language models (LLMs) has opened a new
possibility for educational content authoring: the ability to analyse a
video's transcript and generate contextually appropriate, pedagogically
structured interactive questions
automatically~\cite{brown2020,anthropic2024}. If an instructor can review
and accept AI-generated suggestions rather than composing each question from
scratch, the cognitive overhead of interactive video creation falls
dramatically. The same pipeline that extracts a transcript, classifies its
segments by topic and complexity, and returns a structured list of H5P
interactions can also apply Bloom's taxonomy labels to each suggested
question, helping instructors achieve a deliberate spread of cognitive
difficulty without specialist training in instructional design.

This thesis explores whether combining a purpose-built, user-centred
authoring interface with an AI pipeline based on Anthropic Claude and a
locally hosted Ollama inference server can produce a measurable improvement
in educator productivity and satisfaction when creating H5P interactive
videos.

% ─────────────────────────────────────────────────────────────────────────────
\section{Problem Statement}

The central problem this thesis addresses is the mismatch between the
educational potential of H5P interactive video and the practical usability
of existing authoring environments for university instructors.

The problem manifests through four interacting challenges:

\begin{enumerate}
  \item \textbf{Workflow complexity.} Current WordPress-based H5P workflows
  require educators to navigate at least six distinct screens and understand
  web publishing concepts that are irrelevant to their pedagogical task.
  Task analysis of the existing workflow measured 47 interface interactions
  across six screens to produce one interactive video, with 3.2-second
  median page-load times that interrupt creative flow.

  \item \textbf{High cognitive load.} The extraneous cognitive load imposed
  by the technical interface reduces the mental resources available for
  pedagogical design~\cite{chandler1991,mayer2009}. Instructors report that
  managing the tool consumes more attention than designing the questions
  themselves.

  \item \textbf{Low adoption.} As a direct consequence of the above, adoption
  of H5P interactive video at VNU-HCM has been limited to technically
  proficient faculty. Senior instructors and those in humanities disciplines
  are systematically excluded from using a tool that research demonstrates
  to be educationally effective~\cite{davis1989,hadullo2017}.

  \item \textbf{Absence of AI assistance.} No existing freely available H5P
  authoring tool provides integrated AI assistance for question generation.
  Instructors who could benefit most from AI support---those with limited
  instructional design experience---have no mechanism for obtaining
  context-aware question suggestions tied to the content of their specific
  video.
\end{enumerate}

The \textbf{primary research question} is therefore:

\begin{quote}
  \itshape How can a purpose-built web platform combining user-centred
  interface design with an AI-driven transcript analysis pipeline improve
  the efficiency, quality, and adoption of H5P interactive video authoring
  by university instructors?
\end{quote}

Four supporting research questions guide the investigation:

\begin{enumerate}[label=RQ\arabic*.]
  \item What are the specific usability barriers in current H5P authoring
  workflows that most strongly reduce task completion and satisfaction?
  \item How can user-centred design principles be applied to reduce
  extraneous cognitive load in timestamp-based interactive question
  authoring?
  \item What technical architecture best supports real-time AI-assisted
  question generation within an H5P authoring workflow?
  \item To what extent does the purpose-built platform outperform the
  WordPress-based baseline on measurable usability and productivity metrics?
\end{enumerate}

% ─────────────────────────────────────────────────────────────────────────────
\section{Research Objectives}

The thesis pursues five concrete objectives in order to answer the research
questions above.

\begin{enumerate}
  \item \textbf{Platform development.} Design and implement a full-stack
  web application that enables instructors to upload or link a video, add
  H5P interactive elements at precise timestamps through a visual timeline
  interface, preview the resulting content in real time, and export a
  standards-compliant \texttt{.h5p} package suitable for direct import into
  the VNU-HCM MOOC infrastructure.

  \item \textbf{AI pipeline integration.} Integrate a dual-provider AI
  pipeline that extracts a timestamped transcript from the uploaded video
  (via YouTube automatic captions, uploaded subtitle files, or Whisper audio
  transcription), sends the transcript to a large language model (Anthropic
  Claude or a locally hosted Ollama instance), and returns a validated list
  of H5P interaction suggestions with timestamps, content type, question
  text, and a Bloom's taxonomy classification.

  \item \textbf{User-centred design.} Apply user-centred design methodology
  throughout the interface design process, grounding all design decisions in
  identified instructor needs and validated against usability heuristics.

  \item \textbf{Security implementation.} Implement a security baseline
  aligned to the OWASP Top Ten, covering authentication, authorisation,
  input validation, transport security, and rate limiting.

  \item \textbf{Empirical evaluation.} Conduct a comparative usability study
  measuring task completion rate, time-on-task, error rate, and System
  Usability Scale (SUS) score against the WordPress H5P baseline, and report
  the findings with statistical context.
\end{enumerate}

% ─────────────────────────────────────────────────────────────────────────────
\section{Scope and Limitations}

\subsection{Scope}

The project scope is bounded as follows.

\textbf{Content type scope.} The platform supports the six H5P interaction
types most relevant to interactive video pedagogy: Multiple Choice,
True/False, Fill in the Blanks, Image Hotspot, Drag and Drop, and Mark the
Words. All six are supported for manual authoring; the AI pipeline generates
Multiple Choice, True/False, and Fill in the Blanks suggestions
automatically, as these three types have deterministic schema structures
that can be reliably validated via Zod.

\textbf{Video source scope.} The platform accepts direct video uploads in
MP4, WebM, and MOV formats up to 500~MB, and also accepts YouTube video
URLs. For YouTube-sourced videos, transcripts are extracted directly from
the platform's automatic captions. For uploaded video files, instructors may
upload a WebVTT caption file; alternatively, the system can invoke
\texttt{faster-whisper} locally to perform audio-to-text transcription.

\textbf{User scope.} The primary target user is a university instructor at
VNU-HCM. System administration is handled by a second user role (Admin),
which can manage user accounts, view usage analytics, and export audit logs.

\textbf{Integration scope.} The platform exports standard H5P packages that
can be imported into any LTI-compatible LMS. An LTI~1.3 route is
implemented in the backend for direct LMS integration.

\subsection{Limitations}

\begin{itemize}
  \item \textbf{H5P content types.} The platform does not support the full
  catalogue of 50+ H5P content types. Branching scenarios, Course
  Presentation, and other complex types are out of scope.

  \item \textbf{Evaluation sample.} Usability testing was conducted with
  ten Information Technology students serving as educator proxies, not with
  actual university faculty. While this is a standard and accepted practice
  for early-stage usability evaluation~\cite{nielsen1994}, it introduces a
  limitation in ecological validity.

  \item \textbf{Mobile authoring.} The platform is a responsive web
  application optimised for desktop browsers. Mobile authoring is possible
  but was not a design priority and was not formally evaluated.

  \item \textbf{Concurrent users.} The system was stress-tested at up to
  50 concurrent users. Large-scale institutional deployments serving
  hundreds of simultaneous editors were not evaluated.

  \item \textbf{AI transcription quality.} The quality of AI-generated
  suggestions depends directly on the quality of the transcript. Poor audio
  quality in uploaded videos will degrade Whisper transcription accuracy and,
  consequently, the relevance of AI suggestions.
\end{itemize}

% ─────────────────────────────────────────────────────────────────────────────
\section{Contributions}

This thesis makes the following contributions to educational technology
research and practice.

\begin{enumerate}
  \item \textbf{A reproducible platform architecture.} The full-stack
  architecture combining React~18, Express.js, SQLite/Sequelize, the
  Lumieducation H5P server library, and a dual-provider AI pipeline is
  documented in sufficient detail to enable reproduction and institutional
  adoption.

  \item \textbf{An AI-assisted authoring workflow.} The pipeline that takes
  a raw video transcript through structured LLM prompting to a
  Zod-validated, Bloom's-taxonomy-labelled list of H5P suggestions
  constitutes a novel integration of generative AI into the H5P ecosystem.

  \item \textbf{Prompt engineering for structured H5P output.} The system
  prompts and validation strategy developed to ensure that Claude and Ollama
  return JSON that conforms to H5P content type specifications provide a
  reusable pattern for structured AI output in educational content systems.

  \item \textbf{Empirical usability evidence.} The comparative evaluation
  against the WordPress H5P baseline provides quantitative evidence that a
  purpose-built platform can substantially reduce authoring time and improve
  task completion rates, contributing to the evidence base for user-centred
  design in educational technology.

  \item \textbf{Open-source codebase.} The platform codebase, including all
  AI service logic, prompt definitions, and H5P integration code, is made
  openly available to support further research and institutional adoption.
\end{enumerate}

% ─────────────────────────────────────────────────────────────────────────────
\section{Significance of the Study}

The practical significance of this work is threefold. First, by lowering
the authoring barrier, the platform enables a broader range of VNU-HCM
faculty to produce interactive content and thereby improves the quality
of online learning experienced by thousands of students. Second, the
AI-assisted workflow demonstrates that the gap between AI research and
practical educational tool development can be bridged with modest engineering
effort, encouraging further investment in similar integrations. Third, the
platform aligns directly with Vietnam's national digital transformation
goals~\cite{vietnam_strategy2021} by providing a locally deployable,
open-source solution that does not depend on expensive commercial software.

% ─────────────────────────────────────────────────────────────────────────────
\section{Structure of the Report}

The remainder of this report is organised as follows.

\textbf{Chapter~2 -- Literature Review} surveys the academic and technical
foundations of the work, covering e-learning in Vietnamese higher education,
the H5P framework and its limitations, interactive video research, cognitive
load theory, user-centred design, large language models and prompt
engineering, cloud-native web architecture, and LMS integration standards.
The chapter concludes with an analysis of related systems and identifies the
research gaps that this project addresses.

\textbf{Chapter~3 -- Methodology} describes the research design, the
requirements analysis process, the user-centred design methodology applied
during interface development, the agile development approach, and the
evaluation protocol used to measure usability outcomes.

\textbf{Chapter~4 -- System Design and Implementation} presents the
complete system architecture, database design, backend implementation,
frontend component structure, the AI pipeline in detail, H5P integration
mechanisms, and the security implementation.

\textbf{Chapter~5 -- Results and Evaluation} reports the outcomes of
the comparative usability study, including quantitative task metrics,
SUS scores, and qualitative participant feedback. It also reports on
AI pipeline performance and system responsiveness.

\textbf{Chapter~6 -- Conclusion and Future Work} synthesises the findings,
assesses the contributions in light of the research questions, discusses
limitations, and proposes directions for future research and platform
development.
